% ============================================================
%  Chapter 3: Árboles (Trees)
%  Template: latex-catppuccin-academic  |  Motor: pdflatex
%
%  DEPENDENCIA NUEVA: este capítulo usa el paquete `forest`
%  para la composición tipográfica de árboles. Agregar en
%  config/packages.tex:
%
%      \usepackage[linguistics]{forest}
%
%  El paquete forest requiere tikz, que ya está cargado.
% ============================================================

\chapter{Árboles (\textit{Trees})}
\label{ch:trees}

Los árboles son estructuras de datos jerárquicas y recursivas que resuelven un problema
fundamental de las tablas hash: la ausencia de orden. Como se vio en el capítulo anterior,
un \texttt{HashMap} destruye el orden de los datos para ganar velocidad de acceso; un árbol,
en cambio, preserva ese orden de forma estructural, lo que habilita operaciones como
recorridos ordenados y consultas por rango en $O(\log n)$. Esta propiedad hace de los árboles
la base de estructuras como \texttt{TreeMap} y \texttt{TreeSet} en Java, y de algoritmos
fundamentales en bases de datos e inteligencia artificial.

% ------------------------------------------------------------
\section{Motivación}
\label{sec:trees-motivacion}
% ------------------------------------------------------------

\begin{observacion}{HashMap vs.\ árbol: un intercambio estructural}
Un \texttt{HashMap} aplica una función hash que dispersa las claves de manera uniforme,
destruyendo cualquier relación de orden entre ellas: el acceso es $O(1)$ promedio, pero
preguntas del tipo ``¿cuáles claves están entre $a$ y $b$?'' requieren un recorrido completo
$O(n)$. Un árbol binario de búsqueda preserva el orden de las claves en su estructura,
lo que permite responder esa misma pregunta en $O(\log n)$ mediante un recorrido parcial.
Este intercambio —velocidad puntual versus orden estructural— guía la elección entre ambas
familias de estructuras.
\end{observacion}

\begin{acadtable}[htbp]{tab:hashmap-vs-tree}{Comparación: \texttt{HashMap} vs.\ árbol binario de búsqueda}
    \begin{tabularx}{\textwidth}{l X X}
        \headerrow
        \textbf{Criterio} & \textbf{HashMap} & \textbf{Árbol (BST balanceado)} \\
        \hline
        Búsqueda puntual   & $O(1)$ promedio, $O(n)$ peor caso & $O(\log n)$ \\
        Orden de claves    & No garantizado                     & Sí (orden natural) \\
        Range queries      & $O(n)$                             & $O(\log n)$ \\
        Recorrido ordenado & No                                 & Sí (\textit{in-order}) \\
        Rendimiento predecible & No (depende de colisiones)     & Sí (si está balanceado) \\
    \end{tabularx}
\end{acadtable}

% ------------------------------------------------------------
\section{Terminología Fundamental}
\label{sec:trees-terminologia}
% ------------------------------------------------------------

\begin{definicion}{Árbol}
Un \textbf{árbol} es una estructura de datos jerárquica compuesta por nodos conectados por
aristas, sin ciclos. Formalmente, es un grafo dirigido acíclico con un nodo distinguido
llamado \textbf{raíz} (\textit{root}), desde el cual existe exactamente un camino hacia
cualquier otro nodo. La estructura es inherentemente recursiva: cada nodo es la raíz de
su propio subárbol.
\end{definicion}

El vocabulario básico del que dependen todas las definiciones posteriores es el siguiente:

\begin{itemize}
    \item \textbf{Raíz} (\textit{root}): único nodo sin padre; punto de entrada a toda
          la estructura.
    \item \textbf{Nodo interno} (\textit{internal node}): nodo con al menos un hijo.
    \item \textbf{Hoja} (\textit{leaf}) o nodo externo: nodo sin hijos.
    \item \textbf{Padre} (\textit{parent}) e \textbf{hijo} (\textit{child}): si existe
          una arista de $u$ hacia $v$, entonces $u$ es padre de $v$ y $v$ es hijo de $u$.
    \item \textbf{Hermanos} (\textit{siblings}): nodos que comparten el mismo padre.
    \item \textbf{Ancestro} y \textbf{descendiente}: $u$ es ancestro de $v$ si existe un
          camino de $u$ a $v$; $v$ es descendiente de $u$.
\end{itemize}

Las dos medidas de distancia más importantes —y las que generan más confusión— son la
\textbf{profundidad} y la \textbf{altura}:

\begin{definicion}{Profundidad de un nodo}
La \textbf{profundidad} (\textit{depth}) de un nodo $v$ es la longitud del camino único
desde la raíz hasta $v$, medida en número de aristas. Se define de forma inductiva:
\[
    d(\mathrm{raíz}) = 0, \qquad d(v) = d(\mathrm{parent}(v)) + 1.
\]
La profundidad \emph{sube} desde el nodo hacia la raíz.
\end{definicion}

\begin{definicion}{Altura de un nodo}
La \textbf{altura} (\textit{height}) de un nodo $v$ es la longitud del camino más largo
desde $v$ hasta una hoja en su propio subárbol. Se define recursivamente:
\[
    h(v) =
    \begin{cases}
        0 & \text{si } v \text{ es hoja,} \\[4pt]
        1 + \displaystyle\max_{u \,\in\, \mathrm{hijos}(v)} h(u) & \text{si } v
        \text{ es nodo interno.}
    \end{cases}
\]
La \textbf{altura del árbol} es $h(\mathrm{raíz})$, que coincide con la profundidad
máxima de cualquier hoja. La altura \emph{baja} desde el nodo hacia su hoja más lejana.
\end{definicion}

\begin{observacion}{Dualidad entre profundidad y altura}
Profundidad y altura miden direcciones opuestas a lo largo del mismo eje vertical del árbol.
Para la raíz: $d = 0$ y $h = $ altura total del árbol. Para cualquier hoja: $d = $ su nivel
y $h = 0$. El error más común es calcular la altura de un nodo $v$ como su distancia a la
hoja más profunda del árbol \emph{completo}: esto solo es correcto para la raíz, porque la
raíz contiene todo el árbol como su subárbol. Para cualquier otro nodo, la referencia es
la hoja más lejana \emph{dentro de su propio subárbol}.
\end{observacion}

\begin{figure}[htbp]
    \centering
    \begin{forest}
        for tree={
            circle, draw=CtpOverlay2, fill=CtpSurface0, text=CtpText,
            minimum size=0.95cm, inner sep=1pt,
            l sep=14mm, s sep=7mm,
            font=\sffamily\small
        }
        [\textbf{A}, fill=CtpBlue!35,
            label=right:{\scriptsize $d{=}0,\;h{=}3$}
            [B, label=left:{\scriptsize $d{=}1,\;h{=}2$}
                [D, label=left:{\scriptsize $d{=}2,\;h{=}1$}
                    [H, label=left:{\scriptsize $d{=}3,\;h{=}0$}]
                ]
                [E, label=below:{\scriptsize $d{=}2,\;h{=}0$}]
            ]
            [C, label=right:{\scriptsize $d{=}1,\;h{=}1$}
                [F, label=below:{\scriptsize $d{=}2,\;h{=}0$}]
                [G, label=right:{\scriptsize $d{=}2,\;h{=}0$}]
            ]
        ]
    \end{forest}
    \caption{Árbol con anotaciones de profundidad $d$ y altura $h$ en cada nodo. La raíz
             \textbf{A} tiene $d=0$ y $h=3$; todas las hojas tienen $h=0$. Nótese que $h(D)=1$
             aunque la hoja más profunda del árbol completo (H) está a distancia $1$ de $D$:
             en este caso coincide porque H es hoja del subárbol de $D$.}
    \label{fig:tree-anotado}
\end{figure}

% ------------------------------------------------------------
\section{Subárboles}
\label{sec:subtrees}
% ------------------------------------------------------------

\begin{definicion}{Subárbol}
Dado un nodo $v$ en un árbol $T$, el \textbf{subárbol} de $v$ —denotado $T_v$— es el
subgrafo formado por $v$ mismo, todos sus descendientes, y las aristas que los conectan.
El nodo $v$ actúa como raíz de $T_v$. Formalmente, $T_v$ es el subgrafo inducido por
el conjunto $\{u \in T : v \text{ es ancestro de } u\} \cup \{v\}$.
\end{definicion}

El subárbol es el concepto que hace al árbol una estructura \emph{recursiva}: un árbol
es simplemente la raíz más los subárboles de sus hijos. Esta recursividad es la razón por
la que casi todos los algoritmos sobre árboles —recorridos, búsqueda, inserción,
eliminación— se expresan de forma recursiva de manera natural, y por la que las
demostraciones sobre árboles suelen proceder por inducción sobre la altura.

\begin{ejemplo}{Identificación de subárboles}
En el árbol de la \cref{fig:tree-anotado}, el subárbol del nodo $B$ contiene los nodos
$\{B, D, E, H\}$ con las aristas $B\text{--}D$, $B\text{--}E$ y $D\text{--}H$; su altura
es $2$. El nodo $C$ es raíz de un subárbol distinto $\{C, F, G\}$ con altura $1$. Una
hoja como $H$ constituye un subárbol de tamaño $1$ y altura $0$.
\end{ejemplo}

% ------------------------------------------------------------
\section{Árbol Binario}
\label{sec:binary-tree}
% ------------------------------------------------------------

\begin{definicion}{Árbol binario}
Un \textbf{árbol binario} (\textit{binary tree}) es un árbol en el que cada nodo tiene a
lo sumo dos hijos, denominados \textbf{hijo izquierdo} (\textit{left child}) e \textbf{hijo
derecho} (\textit{right child}). La distinción entre izquierdo y derecho es estructural:
un nodo con solo hijo izquierdo es una configuración distinta de un nodo con solo hijo
derecho, aunque ambos tengan un único descendiente.
\end{definicion}

En la implementación, cada nodo almacena dos referencias, \texttt{left} y \texttt{right},
cualquiera de las cuales puede ser \texttt{null}. Esta representación es la base de los
BST, los \textit{heaps} y los árboles de expresión aritmética.

% ------------------------------------------------------------
\section{Tipos de Árbol Binario}
\label{sec:tree-types}
% ------------------------------------------------------------

Las tres propiedades que se definen a continuación son \textbf{independientes entre sí}:
un árbol puede satisfacer cualquier combinación de ellas. El examen habitualmente presenta
árboles que satisfacen algunas y violan otras, por lo que es fundamental verificar cada
propiedad de forma independiente.

% --- Completo ---
\subsection{Árbol Completo (\textit{Complete Binary Tree})}
\label{subsec:complete}

\begin{definicion}{Árbol binario completo}
Un árbol binario es \textbf{completo} (\textit{complete}) si todos sus niveles están
completamente llenos, excepto posiblemente el último, y el último nivel se llena de
\textbf{izquierda a derecha} sin dejar huecos entre nodos.
\end{definicion}

Esta propiedad garantiza que el árbol puede mapearse biyectivamente a un arreglo de
tamaño $n$ usando la relación $\mathrm{parent}(i) = \lfloor(i-1)/2\rfloor$, sin
desperdiciar posiciones. Es la base estructural de los \textit{heaps} y de la
implementación eficiente de colas de prioridad.

\begin{figure}[htbp]
    \centering
    \begin{minipage}{0.47\textwidth}
        \centering
        \begin{forest}
            for tree={circle, draw=CtpGreen, fill=CtpGreen!20, text=CtpText,
                       minimum size=0.85cm, inner sep=1pt, l sep=11mm, s sep=6mm,
                       font=\sffamily\small}
            [1 [2 [4][5]] [3 [6][,phantom,draw=none,fill=none]]]
        \end{forest}
        \par\smallskip
        {\small\color{CtpGreen}\textbf{Válido} — último nivel lleno de izq.\ a der.}
    \end{minipage}
    \hfill
    \begin{minipage}{0.47\textwidth}
        \centering
        \begin{forest}
            for tree={circle, draw=CtpRed, fill=CtpRed!15, text=CtpText,
                       minimum size=0.85cm, inner sep=1pt, l sep=11mm, s sep=6mm,
                       font=\sffamily\small}
            [1 [2 [,phantom,draw=none,fill=none][5]] [3 [6][7]]]
        \end{forest}
        \par\smallskip
        {\small\color{CtpRed}\textbf{No válido} — hueco en nodo 2 antes de que nodo 3
        tenga hijos.}
    \end{minipage}
    \caption{Árbol completo: válido (izquierda) vs.\ no válido (derecha). La ausencia del
             hijo izquierdo de 2 mientras el nodo 3 ya tiene dos hijos rompe la condición
             de llenado izquierda-derecha.}
    \label{fig:complete-tree}
\end{figure}

% --- Lleno ---
\subsection{Árbol Lleno (\textit{Full Binary Tree})}
\label{subsec:full}

\begin{definicion}{Árbol binario lleno}
Un árbol binario es \textbf{lleno} (\textit{full}) si cada nodo tiene exactamente $0$ o
$2$ hijos. Ningún nodo puede tener exactamente $1$ hijo.
\end{definicion}

Esta propiedad es característica de los \textbf{árboles de expresión aritmética}, donde los
operadores internos tienen siempre dos operandos y los operandos son hojas. Un resultado
combinatorio útil: en un árbol lleno con $n$ nodos internos, el número total de hojas es
exactamente $n + 1$.

\begin{figure}[htbp]
    \centering
    \begin{minipage}{0.47\textwidth}
        \centering
        \begin{forest}
            for tree={circle, draw=CtpGreen, fill=CtpGreen!20, text=CtpText,
                       minimum size=0.85cm, inner sep=1pt, l sep=11mm, s sep=6mm,
                       font=\sffamily\small}
            [1 [2 [4][5]] [3]]
        \end{forest}
        \par\smallskip
        {\small\color{CtpGreen}\textbf{Válido} — todo nodo tiene 0 ó 2 hijos.}
    \end{minipage}
    \hfill
    \begin{minipage}{0.47\textwidth}
        \centering
        \begin{forest}
            for tree={circle, draw=CtpRed, fill=CtpRed!15, text=CtpText,
                       minimum size=0.85cm, inner sep=1pt, l sep=11mm, s sep=6mm,
                       font=\sffamily\small}
            [1 [2 [4][,phantom,draw=none,fill=none]] [3]]
        \end{forest}
        \par\smallskip
        {\small\color{CtpRed}\textbf{No válido} — nodo 2 tiene exactamente 1 hijo.}
    \end{minipage}
    \caption{Árbol lleno: válido (izquierda) vs.\ no válido (derecha). El nodo 2 con un
             solo hijo derecho viola la condición de 0 ó 2 hijos.}
    \label{fig:full-tree}
\end{figure}

% --- Balanceado ---
\subsection{Árbol Balanceado (\textit{Balanced Binary Tree})}
\label{subsec:balanced}

\begin{definicion}{Árbol binario balanceado}
Un árbol binario es \textbf{balanceado} (o equilibrado) si todas sus rutas desde la raíz
hasta una hoja tienen la misma longitud o difieren en a lo sumo $1$:
\[
    \max_{\text{hoja } \ell}\; d(\ell)
    \;-\;
    \min_{\text{hoja } \ell'}\; d(\ell')
    \;\leq\; 1,
\]
donde $d(\cdot)$ denota la profundidad del nodo.
\end{definicion}

La importancia algorítmica del balance es directa: un árbol binario balanceado con $n$
nodos tiene altura $h = O(\log_2 n)$, lo que garantiza $O(\log n)$ en búsqueda, inserción
y eliminación. Si el árbol se desbalancea —degenerando en una lista enlazada— la altura
crece a $O(n)$ y las operaciones pierden toda ventaja sobre una búsqueda lineal.

\begin{figure}[htbp]
    \centering
    \begin{minipage}{0.47\textwidth}
        \centering
        \begin{forest}
            for tree={circle, draw=CtpGreen, fill=CtpGreen!20, text=CtpText,
                       minimum size=0.85cm, inner sep=1pt, l sep=11mm, s sep=6mm,
                       font=\sffamily\small}
            [1 [2 [4][5]] [3 [6][7]]]
        \end{forest}
        \par\smallskip
        {\small\color{CtpGreen}\textbf{Válido} — todas las rutas tienen longitud $2$.}
    \end{minipage}
    \hfill
    \begin{minipage}{0.47\textwidth}
        \centering
        \begin{forest}
            for tree={circle, draw=CtpRed, fill=CtpRed!15, text=CtpText,
                       minimum size=0.85cm, inner sep=1pt, l sep=11mm, s sep=5mm,
                       font=\sffamily\small}
            [1 [2 [4 [8][,phantom,draw=none,fill=none]]
                   [,phantom,draw=none,fill=none]]
               [3]]
        \end{forest}
        \par\smallskip
        {\small\color{CtpRed}\textbf{No válido} — rutas de longitud $3$ y $1$;
        diferencia $= 2 > 1$.}
    \end{minipage}
    \caption{Árbol balanceado: válido (izquierda) vs.\ no válido (derecha). La diferencia
             entre la ruta más larga ($1\to2\to4\to8$, longitud $3$) y la más corta
             ($1\to3$, longitud $1$) supera el umbral de $1$.}
    \label{fig:balanced-tree}
\end{figure}

\begin{observacion}{Independencia de las tres propiedades}
Completo, lleno y balanceado son condiciones ortogonales. Un árbol puede ser completo sin
ser lleno (si el último nivel está parcialmente lleno, algún nodo tendrá $1$ hijo),
lleno sin ser balanceado (todos los nodos tienen $0$ ó $2$ hijos pero las ramas izquierda
y derecha pueden tener longitudes muy dispares), y balanceado sin ser completo (las rutas
pueden diferir en $1$ sin que el llenado sea izquierda-derecha). Para el examen, verificar
cada propiedad de forma independiente y apoyar la respuesta con la evidencia concreta
del nodo que la viola.
\end{observacion}

% ------------------------------------------------------------
\section{Recorridos de un Árbol Binario}
\label{sec:traversals}
% ------------------------------------------------------------

Los tres recorridos clásicos —\textit{pre-order}, \textit{in-order} y \textit{post-order}—
son instancias de la misma plantilla DFS (\textit{depth-first search}) recursiva. La única
diferencia entre ellos es el \textbf{momento en que se procesa el nodo actual} relativo al
recorrido de sus subárboles:

\[
\text{Pre-order:}\;(\mathrm{Raíz},\, L,\, R)
\qquad
\text{In-order:}\;(L,\, \mathrm{Raíz},\, R)
\qquad
\text{Post-order:}\;(L,\, R,\, \mathrm{Raíz})
\]

El prefijo lo indica literalmente: \textit{pre} = raíz primero, \textit{in} = raíz en medio,
\textit{post} = raíz al final. La rama izquierda siempre se visita antes que la derecha;
este orden es invariante en los tres recorridos.

\begin{algorithm}[H]
\caption{Pre-order — (\textit{Raíz, Izquierda, Derecha})}
\label{alg:preorder}
\KwIn{Nodo $v$}
\If{$v = \texttt{null}$}{\Return}
\textbf{procesar}$(v)$\;
\texttt{preOrder}$(v.\mathit{left})$\;
\texttt{preOrder}$(v.\mathit{right})$\;
\end{algorithm}

\begin{algorithm}[H]
\caption{In-order — (\textit{Izquierda, Raíz, Derecha})}
\label{alg:inorder}
\KwIn{Nodo $v$}
\If{$v = \texttt{null}$}{\Return}
\texttt{inOrder}$(v.\mathit{left})$\;
\textbf{procesar}$(v)$\;
\texttt{inOrder}$(v.\mathit{right})$\;
\end{algorithm}

\begin{algorithm}[H]
\caption{Post-order — (\textit{Izquierda, Derecha, Raíz})}
\label{alg:postorder}
\KwIn{Nodo $v$}
\If{$v = \texttt{null}$}{\Return}
\texttt{postOrder}$(v.\mathit{left})$\;
\texttt{postOrder}$(v.\mathit{right})$\;
\textbf{procesar}$(v)$\;
\end{algorithm}

Los tres algoritmos tienen complejidad $\Theta(n)$ en tiempo —cada nodo se visita exactamente
una vez— y $O(h)$ en espacio auxiliar por la pila de recursión, donde $h$ es la altura del
árbol. En el peor caso (árbol degenerado) esto es $O(n)$; en un árbol balanceado, $O(\log n)$.

\begin{acadtable}[htbp]{tab:recorridos}{Comparación de los tres recorridos clásicos}
    \begin{tabularx}{\textwidth}{l X X}
        \headerrow
        \textbf{Recorrido} & \textbf{Orden de visita} & \textbf{Aplicación típica} \\
        \hline
        Pre-order   & Raíz $\to$ Izquierda $\to$ Derecha & Serialización del árbol;
                      reconstrucción de la misma estructura por inserción secuencial \\
        In-order    & Izquierda $\to$ Raíz $\to$ Derecha & Secuencia ordenada en BST;
                      verificación de la propiedad de orden \\
        Post-order  & Izquierda $\to$ Derecha $\to$ Raíz & Liberación de memoria
                      (se procesan hijos antes que el padre); evaluación de
                      árboles de expresión aritmética \\
    \end{tabularx}
\end{acadtable}

\begin{ejemplo}{Recorridos sobre un árbol concreto}
Dado el árbol:

\begin{center}
\begin{forest}
    for tree={
        circle, draw=CtpOverlay2, fill=CtpSurface0, text=CtpText,
        minimum size=0.9cm, inner sep=2pt, l sep=13mm, s sep=9mm,
        font=\sffamily\small
    }
    [A [B [D][E]] [C [,phantom,draw=none,fill=none][F]]]
\end{forest}
\end{center}

\noindent los tres recorridos producen las siguientes secuencias:

\begin{itemize}
    \item \textbf{Pre-order} (Raíz, $L$, $R$): $A,\, B,\, D,\, E,\, C,\, F$
    \item \textbf{In-order} ($L$, Raíz, $R$): $D,\, B,\, E,\, A,\, C,\, F$
    \item \textbf{Post-order} ($L$, $R$, Raíz): $D,\, E,\, B,\, F,\, C,\, A$
\end{itemize}

Nótese el nodo $C$, que tiene únicamente hijo derecho ($F$): en el in-order, al no existir
subárbol izquierdo, se procesa $C$ inmediatamente y luego $F$. Este caso —nodo con un solo
hijo— es el más propenso a errores en examen.
\end{ejemplo}

\begin{observacion}{In-order y BST}
El recorrido in-order de un \textit{Binary Search Tree} (BST) produce las claves en orden
estrictamente ascendente. Esta propiedad no es una coincidencia: es consecuencia directa
de la propiedad de orden del BST, que se estudia en el capítulo siguiente. Verificar que
el in-order produce una secuencia ascendente es la forma más sencilla —y suficiente— de
validar que un árbol es un BST bien construido.
\end{observacion}

% ------------------------------------------------------------
\section{Ejercicios}
\label{sec:trees-ejercicios}
% ------------------------------------------------------------

\begin{ejercicio}
Para el siguiente árbol binario, determine: \textbf{(1)} la altura del árbol, \textbf{(2)}
si es completo, \textbf{(3)} si está lleno, \textbf{(4)} si está balanceado, \textbf{(5)}
la altura y la profundidad de los nodos $1$, $3$, $5$ y $10$, y \textbf{(6)} los recorridos
pre-order, in-order y post-order completos.

\begin{center}
\begin{forest}
    for tree={
        circle, draw=CtpOverlay2, fill=CtpSurface0, text=CtpText,
        minimum size=0.9cm, inner sep=2pt, l sep=13mm, s sep=7mm,
        font=\sffamily\small
    }
    [1
        [2
            [4 [7][8]]
            [5 [9][10]]
        ]
        [3
            [6 [11][12]]
            [,phantom,draw=none,fill=none]
        ]
    ]
\end{forest}
\end{center}
\end{ejercicio}

\begin{ejercicio}
Dado el siguiente árbol:

\begin{center}
\begin{forest}
    for tree={
        circle, draw=CtpOverlay2, fill=CtpSurface0, text=CtpText,
        minimum size=0.9cm, inner sep=2pt, l sep=13mm, s sep=7mm,
        font=\sffamily\small
    }
    [1
        [2
            [4 [,phantom,draw=none,fill=none][8]]
            [5]
        ]
        [3
            [6]
            [7 [10][11]]
        ]
    ]
\end{forest}
\end{center}

\noindent determine las mismas propiedades del ejercicio anterior e identifique con
precisión el nodo (o los nodos) que provocan la violación de cada propiedad que no se
satisfaga.
\end{ejercicio}

\begin{ejerciciocomp}
Implemente en Java los tres recorridos (pre-order, in-order, post-order) de forma recursiva
para un árbol binario genérico. La clase \texttt{Node} debe contener un campo
\texttt{int value} y referencias \texttt{Node left} y \texttt{Node right}. Construya el
árbol del primer ejercicio e imprima los tres recorridos. Verifique sus resultados
comparándolos con los calculados a mano.
\end{ejerciciocomp}


\printbibliography[heading=subbibliography]
